\chapter{Drug Screen (Urine)}
\section*{What Is This Test?}
A urine drug screen (UDS), also known as a urine toxicology screen, is a qualitative or semi-quantitative immunoassay-based test that detects the presence of drugs and/or drug metabolites in urine \cite{moeller2008urine}. The standard urine drug screen panel typically detects: amphetamines, barbiturates, benzodiazepines, cocaine (as benzoylecgonine), marijuana (cannabinoids, THC-COOH metabolite), opiates (morphine, codeine), phencyclidine (PCP), methadone, methaqualone, propoxyphene, and tricyclic antidepressants. Expanded panels include: oxycodone, buprenorphine, fentanyl, synthetic cannabinoids (K2/Spice), MDMA (ecstasy), and gabapentinoids. The immunoassay screen is designed to be highly sensitive (to detect as many true positives as possible) but has limited specificity; confirmation of positive screens by gas chromatography-mass spectrometry (GC-MS) or liquid chromatography-tandem mass spectrometry (LC-MS/MS) is required for definitive identification and is standard of care in forensic, employment, and clinical contexts. Understanding the limitations of the immunoassay is essential: the opiate immunoassay detects natural opiates (morphine, codeine, heroin metabolite 6-MAM) but does NOT reliably detect synthetic/semi-synthetic opioids (oxycodone, hydrocodone, fentanyl, methadone, buprenorphine, tramadol, tapentadol) unless an expanded specific immunoassay is used. The amphetamine immunoassay cross-reacts with many sympathomimetic amines (pseudoephedrine, ephedrine, bupropion, trazodone, selegiline, ranitidine [former]), producing false-positive results. The benzodiazepine immunoassay has poor sensitivity for certain benzodiazepines (lorazepam, clonazepam) because they are excreted primarily as glucuronide conjugates that are not detected unless the urine is pre-treated with glucuronidase.
\section*{Alternative Names}
UDS; Urine Toxicology Screen; Urine Drug Screen; Urine Drugs of Abuse Screen; Toxicology Immunoassay; Drug Test
\section*{Principle}
Competitive immunoassay: drug or drug metabolite in urine competes with a labeled drug-conjugate for antibody binding sites. The more drug present in the urine, the less labeled drug binds to the antibody, producing a signal inversely proportional to drug concentration. Results are qualitative (positive/negative at a defined cutoff concentration). Cutoff concentrations are established by the Substance Abuse and Mental Health Services Administration (SAMHSA) for federal workplace drug testing \cite{samhsa2017mandatory}.
\section*{History}
Urine drug testing began in the 1960s with the development of radioimmunoassays, and in 1972 the enzyme-multiplied immunoassay technique (EMIT) made rapid, inexpensive urine screening widely available. Because immunoassays could not distinguish a drug from cross-reacting compounds, gas chromatography-mass spectrometry (GC-MS) was adopted in the 1970s--1980s as the confirmatory standard. Federal workplace drug testing was formalized by the 1986 Executive Order on a Drug-Free Federal Workplace and the 1988 Mandatory Guidelines, administered by the Substance Abuse and Mental Health Services Administration (SAMHSA) and updated most recently in 2017. Liquid chromatography-tandem mass spectrometry (LC-MS/MS) emerged in the 2000s and extended detection to synthetic opioids and other drugs not covered by older immunoassays. Urine testing has since expanded from workplace and forensic use into routine clinical practice, including pain management monitoring and addiction treatment.
\section*{Why This Test Is Ordered}
\begin{itemize}
  \item Pre-employment, workplace, and federally mandated (DOT) drug testing
  \item Evaluation of suspected drug intoxication or overdose in the emergency department
  \item Monitoring medication compliance in chronic pain management and in addiction treatment programs (e.g., buprenorphine, methadone)
  \item Suspected drug use during pregnancy or by a newborn
  \item Forensic and legal contexts, including post-accident, athletic, and custody testing
  \item Evaluation of altered mental status or unexplained behavior when drug use is suspected
  \item Serial monitoring of inpatients and psychiatric patients for substance use
\end{itemize}
\section*{Primary vs Secondary Test}
The urine drug screen is a primary screening test that is highly sensitive but only presumptive; any positive screen requires secondary confirmatory testing by GC-MS or LC-MS/MS for definitive identification. In chronic pain and addiction monitoring, it functions as a serial surveillance test whose results must be interpreted against the prescribed medication list.
\section*{How This Test Is Performed}
A random urine specimen (approximately 30--45 mL) is collected in a sterile container, with temperature, pH, creatinine, and specific gravity checked for specimen validity (temperature should be 32--38$^{\circ}$C). The sample is screened by automated competitive immunoassay at defined cutoff concentrations, producing qualitative positive or negative results within minutes to hours \cite{rifai2018tietz}. Presumptive positive results are confirmed by gas chromatography-mass spectrometry (GC-MS) or liquid chromatography-tandem mass spectrometry (LC-MS/MS), typically adding 1--3 days. Collection may be directly observed in forensic programs, and chain-of-custody documentation is maintained.
\section*{What Preparation Is Needed}
No special preparation is required; the patient should be adequately hydrated. The collecting program determines whether the collection is observed or unobserved. The patient should be advised which prescription and over-the-counter medications may cause positive results, and should disclose all current medications to the reviewing clinician.
\section*{Contraindications and Precautions}
\begin{itemize}
  \item Prescription and over-the-counter drugs cause cross-reactivity: the amphetamine immunoassay reacts with bupropion, trazodone, selegiline, pseudoephedrine, and labetalol, and the opiate assay reacts with poppy seed ingestion, rifampin, and some quinolones
  \item The opiate immunoassay does not detect synthetic opioids (fentanyl, oxycodone, methadone, buprenorphine) unless an expanded panel is used
  \item Benzodiazepine immunoassays miss lorazepam and clonazepam unless the urine is treated with glucuronidase
  \item Dilution, substitution, and adulteration can produce false-negative results; specimen validity testing detects many of these
  \item A negative result does not exclude use: drugs may be below cutoff, outside the detection window, or not covered by the panel
  \item All positive results should be confirmed by mass spectrometry before any adverse action is taken \cite{standridge2010urine}
\end{itemize}
\section*{Risks}
None. Urine collection is noninvasive and carries no physical risk; possible discomfort or privacy concerns accompany observed collection.
\section*{What Happens After the Test}
A negative result is reported as no drugs detected and requires no further testing. A presumptive positive result is sent for GC-MS or LC-MS/MS confirmation, and in workplace programs a Medical Review Officer (MRO) reviews the result and legitimate prescriptions before it is reported. In clinical settings, confirmed results guide management, including referral for substance use treatment, adjustment of pain medication agreements, or counseling.
\section*{Understanding Results}
\subsection*{Negative}
No drug or metabolite detected above the assay cutoff concentration. Does NOT rule out drug use outside the detection window, dilution (intentional or unintentional), substitution (synthetic urine, other person's urine), or use of drugs not detected by the panel.
\subsection*{Positive (Presumptive Positive)}
Amphetamines: false-positive from bupropion, trazodone, selegiline, pseudoephedrine, and labetalol. Opiates: false-positive from poppy seed ingestion (morphine and codeine), quinolone antibiotics (levofloxacin, ofloxacin), and rifampin. PCP: false-positive from dextromethorphan, ketamine, and venlafaxine. Benzodiazepines: false-negative for lorazepam and clonazepam without glucuronidase treatment. THC: passive inhalation does NOT produce urine THC-COOH levels above the 50 ng/mL cutoff in standard settings. All presumptive positive immunoassay results should be confirmed by GC-MS or LC-MS/MS.
\section*{Normal Results}
Negative for all tested drugs and metabolites. Drug detection windows in urine (approximate): Amphetamines: 2--3 days. Cocaine (benzoylecgonine): 2--4 days (up to 7--10 days with heavy/chronic use). Marijuana (THC-COOH): 3--5 days (single use), 10--15 days (moderate use), 30+ days (chronic heavy use, especially with high body fat percentage). Opiates (morphine/codeine): 2--3 days. Heroin (6-MAM): 6--12 hours (must be collected rapidly). Benzodiazepines: 1--7 days (depends on half-life; short-acting 1--2 days, long-acting diazepam 5--10 days). PCP: 7--14 days with chronic use. Methadone: 2--7 days. Fentanyl: 1--4 days \cite{baselt2017disposition}.
\section*{Follow-up Tests}
Confirmatory testing by GC-MS or LC-MS/MS (gas chromatography-mass spectrometry or liquid chromatography-tandem mass spectrometry). Quantitative levels for specific drugs. Chain-of-custody documentation for forensic/legal specimens. Urine specimen validity testing: creatinine (<20 mg/dL suggests dilution, <2 mg/dL suggests substitution), specific gravity (<1.003 suggests dilution), pH, and oxidant adulterant testing. Alternative matrices: blood (recent use), hair (months of drug use history), oral fluid/saliva (recent use), meconium (neonatal in utero exposure).
\section*{References}
\bibliographystyle{unsrtnat}
\bibliography{references}
\clearpage
\section*{Regulatory Status and Authorization Date}
This chapter describes a test category, not a specific commercial product. FDA, CDSCO, EU IVDR, and MHRA approval or registration dates therefore cannot be assigned without the assay manufacturer, product name, instrument, and intended use. For laboratory-developed tests, an individual agency approval date may not exist. See the Preface for the interpretation of regulatory status.

\clearpage
